\begin{table}[t]
\centering
% \small
\begin{tabular}{p{0.95\columnwidth}}
\toprule
\textbf{Example 1: Code-switching and \textit{imāla}}\\
{{transcription LEB_59}}
\\
\textit{{{translit LEB_59}}}
\\
\midrule
\textbf{Example 2: Register mixing}\\
{{transcription SYR_217}}
\\
\textit{{{translit SYR_217}}}
\\
\bottomrule
\end{tabular}
\caption{
\textbf{Necessity of phonetic transcription.}
Two \textit{LangMedia} utterances (Lebanese, top; Syrian, bottom) show phenomena in \ourbenchmark{} that are visible only in the phonetic tier. In Example 1, the French loan is written as pronounced (\textit{sel danyō}) and \ar{كمان} carries \textit{imāla} (\textit{kamēn}). In Example 2, the same speaker realizes \ar{ق} as \squareipa{q} in the MSA-register adverb \textit{taqrīban} and as \squareipa{ʔ} in the dialectal verb \textit{istaġraʾet}, and keeps the interdental in \textit{ṯmān} while stopping it in \textit{aḫadet}. Arabic script does not represent these effects, even with conventional diacritics, motivating the additional tier.
}
\label{tab:multiglossia}
\end{table}
